\chapter*{Résumé}
\addcontentsline{toc}{chapter}{Résumé}
\argossetchapter{Résumé}


Le transport maritime conteneurisé assure aujourd'hui plus de 80\% du volume du commerce mondial et repose sur un outil en apparence robuste mais structurellement vulnérable : une fois scellé, le conteneur devient un « géant aveugle », dépourvu de visibilité sur l'état de sa cargaison. La majorité des unités sèches, sans alimentation électrique ni connectivité embarquées, ne permettent aucun suivi de leur intégrité physique durant le transit, ce qui engendre des litiges, des avaries non documentées et des pertes financières importantes pour les chargeurs, transporteurs et assureurs. Cet enjeu est renforcé par l'entrée en vigueur, en janvier 2026, des amendements aux conventions SOLAS et MARPOL de l'Organisation Maritime Internationale, qui imposent désormais un signalement rapide de toute perte de conteneur en mer, et prend une acuité particulière en Afrique centrale, où les corridors camerounais (ports de Douala et de Kribi) irriguent également le Tchad et la République Centrafricaine.
Face à ce constat, le projet ARGOS propose de concevoir et de déployer une boîte noire logistique connectée, inspirée des enregistreurs de vol aéronautiques, capable de collecter, horodater et transmettre en temps réel les données d'intégrité physique d'un conteneur (température, humidité, chocs, ouverture non autorisée, position). Le présent mémoire développe cette solution suivant trois volets complémentaires : une étude de marché validant la pertinence commerciale d'un modèle de service indexé sur la valeur de la marchandise, une étude de faisabilité et de viabilité financière démontrant la robustesse économique du projet (valeur actuelle nette positive, taux de rentabilité interne de 89,5 \%), et une étude technique aboutissant à la conception et à la réalisation d'un prototype fonctionnel articulé autour d'un microcontrôleur ESP32, d'un module de communication hybride SIM7000G et d'un récepteur GNSS, couplé à une plateforme web (Flask/PostgreSQL).
ARGOS transforme le conteneur passif en un nœud d'information actif et digne de confiance, ouvrant la voie à une contribution africaine aux enjeux mondiaux de sécurisation du fret maritime.
\vspace{1cm}\noindent 

\textbf{Mots-clés :} traçabilité maritime, boîte noire logistique, Internet des objets (IoT), intégrité des conteneurs, chaîne de blocs (blockchain).

\cleardoublepage
\chapter*{Abstract}
\addcontentsline{toc}{chapter}{Abstract}
\argossetchapter{Abstract}

Containerized maritime transport currently accounts for over 80\% of global trade volume and relies on a tool that is seemingly robust yet structurally vulnerable: once sealed, the container becomes a "blind giant," devoid of any visibility over the condition of its cargo. The majority of dry units, lacking onboard power or connectivity, do not allow for any monitoring of their physical integrity during transit, which leads to disputes, undocumented damage, and significant financial losses for shippers, carriers, and insurers. This issue is heightened by the entry into force, in January 2026, of amendments to the International Maritime Organization's SOLAS and MARPOL conventions, which now mandate the rapid reporting of any container lost at sea, and it takes on particular acuteness in Central Africa, where the Cameroonian corridors (ports of Douala and Kribi) also supply Chad and the Central African Republic.

Faced with this situation, the ARGOS project proposes to design and deploy a connected logistics black box, inspired by aeronautical flight recorders, capable of collecting, time-stamping, and transmitting a container's physical integrity data (temperature, humidity, shocks, unauthorized opening, position) in real time. The present thesis develops this solution across three complementary sections: a market study validating the commercial relevance of a service model indexed to the value of the goods, a feasibility and financial viability study demonstrating the project's economic robustness (positive net present value, internal rate of return of 89.5 \%), and a technical study leading to the design and realization of a functional prototype built around an ESP32 microcontroller, a SIM7000G hybrid communication module, and a GNSS receiver, coupled with a web platform (Flask/PostgreSQL).

ARGOS transforms the passive container into an active and trustworthy information node, paving the way for an African contribution to the global challenges of securing maritime freight.

\vspace{1cm}\noindent 

\textbf{Keywords:} maritime traceability, logistics black box, Internet of Things (IoT), container integrity, blockchain.
